\usetikzlibrary{calc}
\begin{figure}[h]
  \centering
  \resizebox{\linewidth}{!}{%
  \begin{tikzpicture}[node distance=1.2cm and 1.4cm,
    block/.style={
        rectangle,
        draw,
        rounded corners=4pt,
        minimum width=3.2cm,
        minimum height=1.4cm,
        text width=3.0cm,
        align=center,
        font=\small
    },
    arrow/.style={-{Stealth[length=6pt]}, thick},
    label/.style={font=\footnotesize, gray, align=center}
]
 
% Nodes
\node[block, fill=teal!15, draw=teal!60] (trans)
    {\textbf{Transducer}\\[4pt]\footnotesize Ambient\\$\downarrow$\\electrical};
 
\node[block, fill=violet!10, draw=violet!50, right=of trans] (pcc)
    {\textbf{Power conditioning}\\[4pt]\footnotesize Rectifier + PMIC\\impedance matching};
 
\node[block, fill=orange!15, draw=orange!60, right=of pcc] (stor)
    {\textbf{Energy storage}\\[4pt]\footnotesize Capacitor / supercap\\/ rechargeable battery};
 
\node[block, fill=gray!10, draw=gray!50, right=of stor,
      minimum width=1.6cm, text width=1.4cm] (load)
    {\textbf{Load}};
 
% Ambient energy input
\node[label, left=0.8cm of trans] (src) {Ambient\\energy};
\draw[arrow] (src) -- (trans);
 
% Main chain arrows
\draw[arrow] (trans) -- node[above, label]{AC/DC} (pcc);
\draw[arrow] (pcc)   -- (stor);
\draw[arrow] (stor)  -- (load);

  \end{tikzpicture}%
  }
  \caption{Block diagram of a typical energy harvesting system.}
  \label{fig:BD_EHS}
\end{figure}